% !TEX root = ../thesis.tex
% =============================================================================
% Chapter 2: Background and Related Work
% =============================================================================

\chapter{Background and Related Work}
\label{ch:background}

% ---------------------------------------------------------------------------
\section{Metaheuristic Optimisation}
\label{sec:bg-metaheuristics}

A metaheuristic is a high-level, problem-independent strategy for solving optimisation problems of the form
\[
  \min \mathcal{F}(x), \qquad x \in \mathcal{S} \subseteq \mathbb{R}^d,
\]
where $x$ is a $d$-dimensional candidate solution, $\mathcal{S}$ is the bounded feasible search space, and $\mathcal{F}: \mathcal{S} \to \mathbb{R}$ is a black-box objective function that can be evaluated but whose gradient or structure is unknown~\cite{sorensen2013metaheuristics}. Metaheuristics work by iteratively generating, evaluating, and selecting candidate solutions according to heuristic rules that govern how the search space is explored. Unlike exact methods, they make no guarantees of optimality but can find good solutions efficiently across a wide range of problem types.

The field encompasses a broad family of algorithms. Evolution Strategies (ES) maintain a population of candidate solutions and apply Gaussian mutation with self-adaptive step sizes~\cite{back1996evolutionary}. CMA-ES extends this by learning a full covariance matrix that captures the local landscape geometry, and is widely considered the state of the art for continuous black-box optimisation~\cite{hansen2001cmaes}. Differential Evolution (DE) generates trial solutions by combining vector differences sampled from the current population~\cite{storn1997de}.

Despite their success, the design of metaheuristics remains largely a manual endeavour. A vast number of methods have been proposed, many with overlapping mechanisms and limited systematic benchmarking~\cite{aranha2022metaphor}. Choosing the right algorithm for a given problem, configuring its parameters, or designing a new variant requires deep expertise and extensive experimentation. This motivates the pursuit of automated, and increasingly explainable, approaches to algorithm design~\cite{vanStein2025Explainable}.

% ---------------------------------------------------------------------------

\section{Automated Algorithm Design}
\label{sec:bg-aad}

The automated design of optimisation algorithms has evolved through several paradigms, each expanding the scope of what can be automated. The earliest framing is algorithm selection. Rice~\cite{rice1976algorithmselection} formalises the task as learning a mapping from problem features to the best-performing algorithm, drawn from a fixed portfolio of candidates evaluated on a set of problem instances. Meta-learning approaches~\cite{smithmiles2009meta} build on this framework by using landscape features to predict algorithm performance.

Rather than selecting among fixed algorithms, algorithm configuration optimises the hyperparameters of a single algorithm. ParamILS~\cite{hutter2009paramils} applies iterated local search in parameter space, treating algorithm performance on a benchmark as a black-box function of its parameters. These methods can find configurations that substantially outperform defaults, but they operate within a fixed algorithmic structure.

Modular frameworks bridge selection and configuration by decomposing algorithms into interchangeable components. Van Rijn et al.~\cite{vanRijn2016modular} evolve the structure of evolution strategies by treating algorithmic modules (selection, recombination, mutation operators) as parameters of a meta-optimisation problem, and Vermetten et al.~\cite{vermetten2023modde} apply the same principle to DE. Such frameworks can search over millions of algorithm configurations, but they remain limited to the modules their designers included. The design space is large but bounded.

A broader methodology is proposed by Hoos~\cite{hoos2012pbo} under the banner of programming by optimisation. Rather than designing a single algorithm, developers should expose all design decisions as parameters and let automated methods search the resulting space. Genetic Improvement~\cite{petke2018gi} takes this further by applying evolutionary operators directly to source code, modifying and recombining program fragments to improve non-functional properties. These ideas foreshadow the shift to code-generation based approaches, where algorithms are constructed from scratch rather than configured from templates.

% ---------------------------------------------------------------------------
\section{LLMs for Algorithm Design}
\label{sec:bg-llm-aad}

The emergence of LLMs capable of generating complex code has introduced a qualitatively different approach to algorithm design. Rather than searching within a predefined design space, LLMs can generate complete algorithms as executable programs, unconstrained by modular templates. Building on existing taxonomies of LLM--EC integration~\cite{vanStein2024LLaMEA,chauhan2025ecllmsurvey}, we group the relevant literature by the role the LLM plays in the evolutionary loop.

Prompt optimisation applies evolutionary algorithms to optimise the prompts given to LLMs. EvoPrompt~\cite{guo2024evoprompt} uses genetic algorithms to evolve prompts, outperforming human-engineered alternatives. The Automated Prompt Engineer (APE)~\cite{zhou2023ape} takes a similar approach using Monte Carlo search. These methods optimise how we communicate with LLMs but do not use LLMs to generate algorithms directly.

Code generation embeds the LLM itself inside the evolutionary loop. FunSearch~\cite{romeraparedes2024funsearch} demonstrated that LLMs embedded in this way can discover novel mathematical results, generating Python programs for the cap-set problem that surpassed known bounds through a distributed island model with millions of LLM calls. Evolution of Heuristics (EoH)~\cite{liu2024eoh} applies a similar principle to combinatorial optimisation, treating candidate heuristics as individuals in an evolutionary population. Algorithm Evolution using Large Language Model (AEL)~\cite{liu2023ael} is a related approach from the same group. The Self-Taught Optimiser (STOP)~\cite{zelikman2024stop} demonstrates recursive self-improvement where the LLM-generated code is itself a prompt optimiser, creating a meta-learning loop.

Complete metaheuristic generation, the setting this thesis operates in, extends code generation from heuristic components to full algorithms. LLaMEA~\cite{vanStein2024LLaMEA} generates full Python classes with initialisation, state management, and search logic for continuous black-box problems, unlike EoH or AEL which generate small heuristic components for combinatorial problems. The framework combines in-context learning (the LLM sees previous algorithms and their scores) with an evolutionary loop that provides selection pressure. A follow-up work~\cite{vanStein2025Behaviour} extends LLaMEA by logging and analysing the behavioural profiles of generated algorithms, showing that higher-performing algorithms consistently exhibit more intensive exploitation behaviour and faster convergence across six LLaMEA configurations. However, this behavioural information is used only for post-hoc analysis, not as feedback to the LLM.

LLaMEA-SAGE~\cite{vanStein2026LLaMEASAGE} goes a step further by feeding structural information about the generated code back into the prompt. SAGE extracts features from the Abstract Syntax Tree (AST) of each algorithm (graph-theoretic metrics, code complexity measures, and structural statistics) stores them in an archive alongside fitness scores, and uses SHAP analysis~\cite{lundberg2017shap} on a surrogate model to identify which features most influence performance. The resulting guidance is translated into natural-language instructions (e.g., ``try decreasing the cyclomatic complexity'') that augment the mutation prompt, yielding faster convergence on MA-BBOB than vanilla LLaMEA.

Several recent approaches offer alternative search strategies over the same design space. MCTS-AHD~\cite{zheng2025mctsahd} replaces the evolutionary loop with Monte Carlo Tree Search, organising generated heuristics in a tree structure that balances exploration and exploitation. LLM-driven Heuristic Neighbourhood Search (LHNS)~\cite{xie2025lhns} abandons population-based evolution in favour of single-solution neighbourhood search with ruin-and-recreate operators. Both achieve competitive results on combinatorial benchmarks, demonstrating that the search strategy over LLM-generated code is itself an active research question.

% ---------------------------------------------------------------------------
\section{Benchmarking and Evaluation}
\label{sec:bg-benchmarking}

Rigorous benchmarking is essential for comparing generated algorithms and providing meaningful feedback to the LLM. The Black-Box Optimisation Benchmark (BBOB)~\cite{hansen2021coco} defines 24 noiseless test functions spanning five function groups (separable, moderate conditioning, high conditioning, multimodal with adequate structure, multimodal with weak structure). Many-Affine BBOB (MA-BBOB)~\cite{vermetten2023mabbob} extends this by constructing novel instances through affine combinations of BBOB functions, providing much greater instance diversity while preserving known landscape properties. MA-BBOB is the benchmark used in recent LLaMEA versions~\cite{vanStein2026LLaMEASAGE} and in this thesis.

The evaluation infrastructure underneath these benchmarks is IOHexperimenter~\cite{deNobel2024iohexperimenter}, which runs algorithms on benchmark functions, enforces budget limits, and logs full trajectories; IOHanalyzer~\cite{wang2022iohanalyzer} provides statistical analysis and visualisation of the logged data. BLADE~\cite{vanStein2025BLADE} builds on IOHexperimenter to provide a standardised evaluation framework specifically for LLM-driven algorithm design, including prompt management, candidate logging, and behavioural metric computation.

The primary performance signal used in this thesis is AOCC~\cite{vanStein2024LLaMEA}, an anytime performance metric that captures not just the final solution quality but the entire convergence trajectory. Unlike fixed-target metrics that measure time-to-threshold, AOCC rewards algorithms that converge quickly \emph{and} reach good final solutions. It is the standard performance metric in the BLADE framework (see \cref{eq:aocc} for the formal definition).

% ---------------------------------------------------------------------------
\section{Trajectory-Based Analysis of Metaheuristics}
\label{sec:bg-behaviour}

A growing body of work analyses algorithm behaviour through the lens of its search trajectory. Search Trajectory Networks (STNs)~\cite{ochoa2021stn} discretise the search space and represent trajectories as graphs, enabling visual comparison of algorithm dynamics. DynamoRep~\cite{cenikj2023dynamorep} uses trajectory-based population dynamics, specifically per-generation statistics of objective and position distributions, to classify optimisation problems. Kostovska et al.~\cite{kostovska2022trajectory} combine landscape features with algorithm-internal state variables (step size, covariance matrix) computed along trajectories for warm-started algorithm selection. Cenikj et al.~\cite{cenikj2026features} survey the broader meta-feature landscape spanning problem, trajectory, and algorithm-internal signals. Across this body of work, however, trajectory-based features have been used for algorithm selection and problem classification, not as feedback signals within an algorithm design loop.
